% ======================================================================================
% Nom     : reports/latex/sections/01_introduction.tex
% Rôle    : Section d'introduction du rapport.
% Auteur  : Maxime BRONNY
% Version : V1
% Cadre   : Réalisé dans le cadre du cours Techniques d'apprentissage artificiel
%           Master 1 Informatique Big Data
% Usage   : Fichier inclus dans main.tex lors de la compilation du rapport.
% ======================================================================================

\section{Introduction}
\label{sec:introduction}

L'évaluation du risque de crédit est l'une des activités les plus sensibles du
secteur bancaire : chaque prêt accordé repose sur un pari, celui que l'emprunteur
remboursera. Historiquement confiée à des analystes humains s'appuyant sur des
grilles de notation, cette évaluation est aujourd'hui largement automatisée grâce
à l'apprentissage automatique, qui permet d'exploiter les données historiques de
remboursement pour estimer la probabilité de défaut d'un nouveau client. Le
\emph{credit scoring} constitue ainsi l'un des cas d'usage les plus anciens et les
mieux documentés de la classification supervisée \cite{hosmer2000logistic}.

Ce projet, réalisé dans le cadre du cours de Techniques d'Apprentissage
Artificiel, a pour objectif de construire un pipeline complet d'apprentissage
automatique appliqué à la prédiction du risque de crédit, puis de comparer
rigoureusement trois méthodes de classification de familles différentes : la
régression logistique, l'arbre de décision CART et les $k$ plus proches voisins
($k$-NN). Un parti pris fort structure l'ensemble du travail : les trois
algorithmes, ainsi que toutes les métriques d'évaluation, sont implémentés
\emph{from scratch} en Python~\cite{python_psf}, à l'aide de la seule
bibliothèque de calcul numérique NumPy \cite{harris2020numpy}. La bibliothèque scikit-learn
\cite{pedregosa2011sklearn} n'intervient qu'en fin de chaîne, uniquement comme
référence permettant de \emph{vérifier} la justesse de nos implémentations. Elle
n'est jamais utilisée comme solution.

Précisons un point d'historique : la première version de ce projet avait été
développée en langage C, avec un découpage train/test simple et deux modèles.
Le projet a été entièrement réécrit en Python pour répondre aux consignes de la
matière ; cette réécriture a été l'occasion d'enrichir le protocole expérimental
(ajout d'un ensemble de validation dédié au choix des hyperparamètres, ajout
d'une troisième méthode) et de valider systématiquement les implémentations par
comparaison avec scikit-learn. L'ancienne version est conservée à titre d'archive
mais n'intervient plus dans le présent travail.

Ce rapport est organisé comme suit. La section~\ref{sec:contexte} présente le
contexte métier et la problématique. La section~\ref{sec:donnees} décrit le jeu
de données utilisé et la section~\ref{sec:pretraitement} les étapes de
prétraitement, dont le découpage entraînement/validation/test. La
section~\ref{sec:methodologie} détaille le protocole expérimental et les
métriques retenues, puis la section~\ref{sec:algorithmes} expose le principe et
l'implémentation de chacun des trois algorithmes. Les résultats sont présentés
en section~\ref{sec:resultats}, analysés de manière critique en
section~\ref{sec:analyse}, et la section~\ref{sec:limites} discute les limites
du travail et les améliorations envisageables, avant la conclusion.
